\chapter*{Introduction}

This document is at the same time the documentation and an example of
\texttt{slimthesis}, a \LaTeX{} style for bachelor's and master's theses
supervised at the SLIM laboratory of the Università degli Studi di Milano.
Its source is the starting point for a new thesis: replace the text of the
chapters with your own, keep the structure, and the style takes care of the
cover page, the typography and the bibliography.

All the layout is in a single file, \texttt{slimthesis.sty}. The document
itself (\texttt{main.tex} and the files it includes) only contains the thesis
data and the text, so that the style can be updated without touching your
work, and your work never needs to touch the style.

\paragraph{Requirements}

The style relies on OpenType fonts, therefore it must be compiled with
\emph{LuaLaTeX} or \emph{XeLaTeX}; pdf\LaTeX{} is not supported. Besides a
recent \TeX{} Live distribution, you need:
\begin{itemize}
  \item the \emph{STIX Two Math} font, used for mathematics (the Fedora
    package is \nolinkurl{stix-fonts});
  \item the \emph{Libertinus} fonts, used for text (the
    Fedora package is \nolinkurl{texlive-libertinus-fonts});
  \item the \emph{Figtree} font, used for headings, which is shipped with the
    template in \nolinkurl{fonts/Figtree/};
  \item a few \LaTeX{} packages that minimal installations may lack, such as
    \texttt{apptools}, \texttt{ocgx2} and \texttt{pdfx} (on Fedora:
    \nolinkurl{texlive-apptools}, \nolinkurl{texlive-ocgx2} and
    \nolinkurl{texlive-pdfx});
  \item \emph{biber}, which processes the bibliography.
\end{itemize}

\paragraph{Compiling}

A complete build runs LuaLaTeX, then biber, then LuaLaTeX twice more, so
that the bibliography, the table of contents and all cross-references are
resolved:
\begin{Verbatim}
lualatex main && biber main && lualatex main && lualatex main
\end{Verbatim}
Alternatively, \texttt{latexmk} runs all the required steps automatically;
the file \texttt{latexmkrc} shipped with the template tells it to use
LuaLaTeX and biber, so that it is enough to type
\begin{Verbatim}
latexmk main
\end{Verbatim}
and \verb|latexmk -c| removes the auxiliary files. XeLaTeX can be used in
the same way, replacing \texttt{lualatex} with \texttt{xelatex}.

\paragraph{PDF/A}

The thesis is produced as a PDF/A-2b file, the archival format usually
required when the final version is deposited, through the \texttt{pdfx}
package. The title, the author, the degree course and the keywords given
with \cs{keywords} are stored in the metadata of the file; the style writes
them in the file \texttt{main.xmpdata}, read by \texttt{pdfx}, which is
regenerated at each compilation and should not be edited. Conformance can be
checked with the free validator veraPDF (\url{https://verapdf.org}). The
option \texttt{pdfa=false} produces an ordinary PDF file instead
(Section~\ref{sec:programme}).

\paragraph{Files}

\begin{description}
  \item[\texttt{main.tex}] the main file: programme, thesis data and the list
    of chapters;
  \item[\texttt{slimthesis.sty}] the style; normally it does not need to be
    modified;
  \item[\texttt{frontmatter/}] dedication, epigraph and acknowledgements;
  \item[\texttt{chapters/}] one file per chapter;
  \item[\texttt{bibliography.bib}] the bibliographic references;
  \item[\texttt{style/}] your own macros (\texttt{symbols.tex}) and the
    colours used by the figures (\texttt{colors.tex});
  \item[\texttt{images/}] figures, the logos used on the cover
    (\texttt{images/coverPage/}) and the Creative Commons buttons
    (\texttt{images/license/});
  \item[\texttt{fonts/}] the Figtree font.
\end{description}

Chapter~\ref{ch1} explains how to set up the thesis, Chapter~\ref{ch2} how
to write it, Chapter~\ref{ch3} how to cite and link, and the last chapter
lists a few checks to do before submitting.
